\section{Limitations}
\label{sec:limitations}

We explicitly delineate several structural and methodological limitations of this study:

\begin{enumerate}
    \item \textbf{Null-Baseline Dependence}: Both utility $U(S)$ and behavioral reliance $R(S)$ are measured relative to a chosen unimodal null baseline $B_{\text{null}}$. While our primary evaluations utilize the \emph{removed} condition (true unimodal ablation), alternative nulls (such as black, mean, or noise images) yield different baseline accuracies and can shift both $S^*$ and $\mi$.
    
    \item \textbf{Conflation of Gate Sharpness and Crossing Timing in $\mi$}: The Misalignment Index integrates normalized reliance over $[S^*, 1]$. Consequently, a model that crosses zero utility very early (small $S^*$) with a gradual gating transition can yield a similar $\mi$ value to a model that crosses late with sharp, persistent reliance.
    
    \item \textbf{Residual Reliance in Smooth Sigmoid Gating}: In parametric gating models utilizing smooth activations (such as sigmoid or softmax), the gate weight assigned to a degraded modality rarely reaches strict mathematical zero, leading to non-zero residual reliance even under complete informational corruption.
    
    \item \textbf{4-Bit Quantization Effects in LLaVA-1.5-7B}: Due to GPU accelerator memory constraints on single T4 instances, LLaVA-1.5-7B is evaluated in 4-bit NormalFloat (NF4) quantization. Quantization may alter predictive logit calibration relative to full-precision weights.
    
    \item \textbf{Evaluation Benchmark Scope}: Our empirical VLM evaluation is conducted exclusively on the 4-choice multiple-choice formulation of A-OKVQA ($N=500$ validation items), rather than open-ended free-form generation.
    
    \item \textbf{Prompt-Format-Dependent Letter Mass}: In modern instruction-tuned VLMs, the proportion of softmax probability mass assigned to choice letters at the initial generation position is highly sensitive to the prompt template and assistant pre-fill prefix. Models whose letter mass fails the $\ge 0.90$ threshold cannot be reliably evaluated under a fixed single-token multiple-choice readout.
    
    \item \textbf{Exclusion of InternVL2}: InternVL2-2B was dropped from our multi-model benchmark due to an unresolvable library incompatibility between its custom remote repository code and `transformers>=4.50` (missing `all_tied_weights_keys` attribute during model instantiation).
    
    \item \textbf{Autonomous Driving (BEV) Study --- Proposed, Not Frozen}: A standalone Camera-LiDAR BEV architecture and physical corruption suite were developed as a feasibility probe; however, no pre-declared hypotheses were formalized for this case study. Per the pre-registration protocol, the BEV analysis remains outside the frozen paper scope and is deferred to future work pending formal hypothesis declaration.
\end{enumerate}

\begin{itemize}
\item \textbf{Null-baseline dependence.} Conclusions for individual
  models depend on whether the no-input baseline is a removed image or a
  black image; the divergence between the two is about two orders of
  magnitude larger for SmolVLM and LLaVA-1.5 than for Qwen2-VL, and part of
  their residual reliance may reflect the presence of image tokens.
\item \textbf{Multiple comparisons.} The criterion is evaluated in 42
  combinations without correction, and the two that satisfy it overlap.
\item \textbf{Operationalization after the fact.} The interval-based form
  of the pre-registered criterion was fixed after seeing data.
\item \textbf{Normalization by $R(0)$.} $\hat R$ and MI are unreliable
  when clean-input reliance is small, as in the fusion study.
\item \textbf{Misalignment Index.} Undefined in most real conditions;
  validated mainly on synthetic models; conflates gate sharpness and
  timing.
\item \textbf{Scope.} One dataset, one multiple-choice format, three
  models (one in 4-bit), one swap pairing per sample, and no autonomous
  driving evaluation.
\end{itemize}
